\chapter{Quadratic curvature energy under varifold convergence}\label{sec:willmore-varifold-energy}
The first-variation convention is Definition \ref{def:first-variation}, based on \cite{Sim18}, Chapter 8, equations (2.3)--(2.10), pp.~239--240. A singular boundary part must not be discarded when discussing continuity of curvature energy. We formulate an extended energy through smooth ambient tests and prove its lower semicontinuity using the representation results already established in the measure chapter.

\begin{definition}[An extended quadratic first-variation energy]\label{def:varifold-quadratic-energy}
For an $m$-varifold $V$ in an open set $U\subset\mathbb R^N$, with weight $\mu=\|V\|$, define
\[
 \mathcal E_U(V)=\sup_{X\in C_c^1(U;\mathbb R^N)}
 \left\{\delta V(X)-\int_U|X|^2d\mu\right\}\in[0,+\infty].
\]
The zero test makes the supremum nonnegative. This definition uses the entire first variation. In dimension two it agrees with one quarter of the squared generalized mean curvature when that curvature represents the first variation without a singular part and is globally square integrable, as proved next.
\end{definition}

\begin{theorem}[Finite quadratic energy is equivalent to square-integrable mean curvature]\label{thm:varifold-energy-dual-representation}
The quantity $\mathcal E_U(V)$ is finite if and only if there is $H\in L^2(\mu;\mathbb R^N)$ such that
\[
 \delta V(X)=-\int_UH\cdot X\,d\mu
 \quad\text{for every }X\in C_c^1(U;\mathbb R^N).
\]
In that case $H$ is unique up to $\mu$-null sets and
\[
 \mathcal E_U(V)=\frac14\int_U|H|^2d\mu.
\]
The first variation therefore has no singular boundary part in $U$. This is a dual formulation of the mean-curvature representation in \cite{Sim18}, pp.~239--240; the $L^2$ equivalence is proved here.
\end{theorem}
\begin{proof}
Suppose $E=\mathcal E_U(V)<\infty$. For a fixed test $X$, applying the definition to $tX$ for every real $t$ gives
\[
 t\delta V(X)-t^2\|X\|_{L^2(\mu)}^2\leq E.
\]
If $\|X\|_2=0$, this forces $\delta V(X)=0$. Otherwise optimize the quadratic polynomial in $t$ to obtain
$|\delta V(X)|\leq2\sqrt E\|X\|_2$.
Thus the first variation is well defined and bounded on the $L^2$ equivalence classes of smooth compactly supported fields.

These fields are dense in $L^2(\mu;\mathbb R^N)$. Indeed Theorem \ref{thm:radon-lp-density} approximates each component by a continuous compactly supported function. Extending such a function by zero and mollifying gives uniform smooth approximation with support in a fixed compact neighborhood inside $U$; local finiteness of $\mu$ converts uniform convergence to $L^2$ convergence. The bounded functional extends to the Hilbert space. Lemma \ref{lem:hilbert-functional-representation} gives a unique vector field $-H$ representing it, with $\|H\|_2\leq2\sqrt E$.

Conversely, if such $H$ exists, completion of the square gives
\[
 \delta V(X)-\int|X|^2d\mu
 =\frac14\int|H|^2d\mu-\int|X+H/2|^2d\mu.
\]
The right side is bounded above by the claimed energy. Smooth compactly supported approximations to $-H/2$ in $L^2$ make the last term tend to zero, proving equality. The same density proves uniqueness. The representing vector measure $-H\mu$ is locally finite by Cauchy--Schwarz and local finiteness of $\mu$, and is absolutely continuous with respect to $\mu$, so there is no additional singular first-variation measure.
\end{proof}

\begin{theorem}[Lower semicontinuity of quadratic varifold energy]\label{thm:varifold-quadratic-energy-lsc}
If $V_j\rightharpoonup V$ as varifolds in $U$, then
\[
 \mathcal E_U(V)\leq\liminf_{j\to\infty}\mathcal E_U(V_j).
\]
In particular, if the liminf is finite, the limit has a generalized mean curvature in $L^2(\|V\|)$ representing its entire first variation, with the corresponding squared-norm bound. This argument uses the first-variation formula of \cite{Sim18}, equation (2.3), and does not assume integrality.
\end{theorem}
\begin{proof}
For a fixed $X\in C_c^1(U;\mathbb R^N)$, both $(x,S)\mapsto\operatorname{div}_SX(x)$ and $(x,S)\mapsto|X(x)|^2$ are continuous with compact support on the Grassmann bundle. Thus
\[
 \delta V_j(X)-\int|X|^2d\|V_j\|
 \longrightarrow \delta V(X)-\int|X|^2d\|V\|.
\]
Each term before the limit is bounded above by $\mathcal E_U(V_j)$. For every fixed test the limiting expression is therefore at most the liminf of the energies. Taking the supremum over tests proves the inequality, including an infinite liminf. If it is finite, apply Theorem \ref{thm:varifold-energy-dual-representation} to identify the limit curvature and its norm. The proof does not interchange a supremum with a limit as an equality.
\end{proof}

\begin{proposition}[Comparison with the energy of a smooth immersion]\label{prop:smooth-varifold-willmore-comparison}
Let $F:M^2\to\mathbb R^N$ be a smooth closed immersion, and let $V_F$ be its induced varifold with parameter multiplicity. Then
\[
 \mathcal E_{\mathbb R^N}(V_F)\leq\mathcal W(F).
\]
For an embedding there is equality. For every compact set $K\subset U$ and any varifold of finite quadratic energy,
\[
 \|\delta V\|(K)\leq2\sqrt{\mathcal E_U(V)}\sqrt{\|V\|(K)}.
\]
\end{proposition}
\begin{proof}
The smooth first variation gives $\delta V_F(X)=-\int_M\vec H_F\cdot X(F)d\mu_M$, and the weight identity gives $\int|X|^2d\|V_F\|=\int_M|X(F)|^2d\mu_M$. Completing the square on the parameter surface bounds their difference by $\frac14\int_M|\vec H_F|^2$. Take the supremum. For an embedding, the geometric mean curvature defines a vector field on the image representing the entire first variation, so Theorem \ref{thm:varifold-energy-dual-representation} gives equality. The weaker assertion for an arbitrary immersion suffices here and avoids requiring a separate identification of curvatures on coincident sheets.

For the last estimate, finite energy represents the first variation by $-H\|V\|$. Its total variation on $K$ is $\int_K|H|d\|V\|$. Cauchy--Schwarz bounds this by $\sqrt{\|V\|(K)}\|H\|_{L^2(U)}$, and the norm identity gives the claimed constant. The same argument applies to relatively compact Borel sets.
\end{proof}

\begin{corollary}[An integral compactness consequence, with its external input]\label{cor:willmore-integral-compactness-conditional}
Let $V_j$ be integral two-varifolds in $U$ with locally uniformly bounded weights and $\sup_j\mathcal E_U(V_j)<\infty$. Then a subsequence converges to an integral two-varifold $V$, and
$\mathcal E_U(V)\leq\liminf_j\mathcal E_U(V_j)$ after taking a subsequence attaining the energy liminf. The integrality conclusion uses the externally proved Allard compactness theorem \ref{thm:allard-compactness}; it is not a consequence of the Hilbert-space duality alone.
\end{corollary}
\begin{proof}
First pass to a subsequence whose energies tend to the liminf. Proposition \ref{prop:smooth-varifold-willmore-comparison} bounds the total variation of $\delta V_j$ uniformly on each relatively compact set from the assumed weight and energy bounds. Apply Theorem \ref{thm:allard-compactness}, whose hypotheses are now verified, to obtain an integral varifold subsequential limit. Theorem \ref{thm:varifold-quadratic-energy-lsc} supplies the energy inequality. This proves the deduction conditional on the explicitly named external compactness theorem; it does not replace its geometric proof.
\end{proof}

\begin{corollary}[A bounded-energy sequence of closed surfaces has local compactness]\label{cor:closed-surfaces-willmore-compactness}
Let $F_j:M_j^2\to\mathbb R^N$ be smooth closed embedded surfaces with $\mathcal W(F_j)\leq\Lambda$. Their induced varifolds have a subsequence converging locally to an integral two-varifold $V$ with $\mathcal E_{\mathbb R^N}(V)\leq\Lambda$. Its weight satisfies
\[
 \|V\|(B_r(y))\leq\Lambda r^2
\]
for every ball, and has no atoms. The limit is allowed to be zero. Integrality again uses Theorem \ref{thm:allard-compactness} externally.
\end{corollary}
\begin{proof}
Theorem \ref{thm:willmore-uniform-quadratic-area} gives uniform mass bounds on every ball and hence on every compact set. Proposition \ref{prop:smooth-varifold-willmore-comparison} identifies the energy of each embedding with its classical Willmore energy. Corollary \ref{cor:willmore-integral-compactness-conditional} gives the claimed limit and energy bound. Weak varifold convergence implies weak convergence of the weights, since the Grassmannian is compact. The open-set lower semicontinuity of Radon measures then yields
$\|V\|(B_r(y))\leq\liminf_j\|V_j\|(B_r(y))\leq\Lambda r^2$.
Let $r\downarrow0$ to get $\|V\|(\{y\})=0$. No fixed area lower bound, tightness at infinity or nonvanishing limit was assumed.
\end{proof}
